% Authors: Keerthi Rapolu & Sreeja Katta

\subsection{Anomaly Root Cause Analysis and Prevention Feedback}

The RCA module converts recurring divergence patterns into policy suggestions,
closing the governance loop between runtime monitoring and future pre-billing
intervention. After each generation, \texttt{RootCauseAnalyzer} queries
historical incidents grouped by workload type and incident category ---
over-provisioned requests, idle cluster patterns, runaway execution, and
undeclared token budget --- and synthesizes a \texttt{Policy} object when the
same pattern recurs at least three times within a lookback window. Confidence
scales with pattern frequency up to a ceiling of 0.95, and only policies above
the 0.60 threshold enter the registry.

This mechanism addresses a gap that static pre-billing policy cannot fill.
Preconfigured rules handle known divergence patterns, but novel or low-frequency
patterns are not covered at system initialization. The feedback loop fills that
gap incrementally: workloads that escape pre-billing intervention, behave
anomalously at runtime, and are flagged by IFS contribute to learned policies
that catch the same pattern in subsequent submissions. The Exp~6 convergence
result is partly attributable to this accumulation --- Full PBCP substantially
outperforms the No Phase~3 ablation because learned policies progressively
reduce the baseline rate of uncaught divergence.

Two limitations apply to the current implementation. First, the feedback loop
operates on synthetic incident data, so confidence and frequency estimates may
not reflect the noisier incident streams in real enterprise environments. Second,
the pattern-counting logic does not weight incidents by recency or cost severity,
meaning a cluster of low-cost incidents and a smaller cluster of high-cost
incidents are treated identically. Both limitations are addressable in a
production deployment and represent near-term extensions rather than design
defects.
